\section{死刑复核程序}
\begin{enumerate}
  \item 基本原则。
        \begin{enumerate}
          \item \bs{未复核判决不生效}。死刑复核程序是死刑案件的终审程序。
          \item 裁判做出后，生效前核准。
          \item 程序启动\zd{自动}性。
          \item 逐级报请，再审可以越级申诉。
          \item 二审 + 复核，可不再另行复核。\zd{复核主体和二审主体同一}。
        \end{enumerate}
  \item 当事人。\zx{辩护权未消失}。\al{核准死刑后，甲不得再为乙进行辩护 \cuo}。
  \item 检察院。最高法裁判文书下发前，最高检\zd{可以}向最高法\zd{提出意见}。
  \item 法院。
        \begin{enumerate}
          \item 死立执。最高法核准。
                \begin{enumerate}
                  \item \yi 不上不抗，上诉期满 10 日。
                        \er 上诉/抗诉 \yjt 触发二审 \yjt 维持死/改活 \yjt 上报/终审。
                        % https://q.uiver.app/?q=WzAsMyxbMCwyLCLmrbvliJHvvIjkuK3pmaLvvIkiXSxbMCwxLCLlj6/ku6XlkKblhrMiXSxbMCwwLCLmnIDpq5jpmaIiXSxbMCwxLCLkuIrmiqUiLDFdLFsxLDIsIuS4iuaKpSIsMV0sWzAsMSwi5o+Q5a6hIiwyLHsiY3VydmUiOjN9XSxbMCwxLCLlj5Hlm54iLDAseyJjdXJ2ZSI6LTN9XV0=
\[
  \begin{tikzcd}[row sep=huge]
    \txt{最高院} \\
    \txt{可以否决} \\
    \txt{死刑（中院）}
    \arrow["\txt{上报}"{description}, from=3-1, to=2-1]
    \arrow["\txt{上报}"{description}, from=2-1, to=1-1]
    \arrow["\txt{提审}"', shift right=8, curve={height=6pt}, from=3-1, to=2-1]
    \arrow["\txt{发回}", shift left=8, curve={height=-6pt}, tail reversed, no head, from=3-1, to=2-1]
  \end{tikzcd}
\]

                  \item \zx{同案犯死刑没上诉，其他人上诉，一锅端，带入二审}。附民带不动刑事。
                  \item 复核程序。\bs{讯问、被动听、通报、审共犯}。
                        \begin{enumerate}
                          \item \zx{复核法官仅 3 人、书面审理}（阅卷 + 提讯）。\al{开庭审理 \cuo 书面审理 \dui}。
                          \item 高院、最高法\bs{应当讯问}被告人。
                          \item 被动听。该辩的法庭上都辩完了。\al{应当听取意见 \cuo。辩护律师提出要求的，应当听 \dui}。
                          \item 应当将死刑复核结果\bs{通报最高检}。
                          \item \bs{共犯一并审查}。
                        \end{enumerate}
                  \item 处理。\zx{维、改、发}。
                        \begin{enumerate}
                          \item \zd{没错核准}。原判事实、法律正确、量刑适当、程序合法。
                          \item \zd{瑕疵纠正}。某一事实或法律引用不完全准确、规范。\zd{纠正后核准} \yjt 继续执行。
                          \item \zd{依法不应当判处死刑}。\bdy 要判成死缓。\bs{复核处理里唯一的可改可发}。\al{怀孕}。\al{将甲改成死缓 \cuo}。
                          \item \bs{见错就发}。\bs{不分人不分罪}，全案\zd{裁定}不予核准，发回重审。\zd{可不开庭做决定，径行发回}。
                                \begin{itemize}
                                  \item \zx{前提不是加重罚}！！\al{中院判甲死缓，省高院认为应判处甲死缓 + 限制缓刑 \yjt 全案发回 \cuo}。
                                  \item \zd{事实}不清、\zd{证据}不足。唯一的不能改。
                                  \item 漏罪、新罪需要追诉。\al{确认甲实施了其他犯罪，应裁定继续执行死刑 \cuo}。
                                  \item 出现\zd{新的}影响定罪和量刑的事实、证据。\al{重大立功、另有犯罪}。
                                  \item 原审违反法定\zd{程序}。
                                \end{itemize}
                        \end{enumerate}
                  \item \zx{死刑复核的发回重审}。发回一审 \hz 二审。
                        \begin{enumerate}
                          \item \zd{一般不得发回一审}。例外。\bs{发回的高院是复核法院，复核法院可以发回重审 \hz 提审}（\zd{没实质审理过}）。
                                \al{一审 A 中，二审 B 高，最高法可发回 A 中 \hz B 高 \dui}。
                          \item \zd{原则重组合议庭}。例外。\yi 新的事实、证据。\er \zd{量刑有偏差}（不应判死刑）。\san 原合议庭无错。
                                \al{应另行组成合议庭 \cuo}。
                          \item \zx{这里考题可能还会涉及合议庭人数问题，要看是用几审程序审的}。
                          \item 若发回一审重审的，应当开庭。
                          \item 若发回二审重审的，可以直接改判。\bs{事实、证据、纠程序，必开庭}。
                          \item 二审带复核可以直接改判，所作判决为终审判决。
                        \end{enumerate}
                \end{enumerate}
          \item 死缓。高院核准。
                \begin{enumerate}
                  \item 没错。二审/复核 \yjt 维持/核准。
                  \item 法律错、量刑错/量刑过重（死缓复核） \yjt 应当改判。
                  \item 事实不清、证据不足/新的事实证据（死缓复核） \yjt 发回 \hz 改判。
                  \item 程序违法 \yjt 应当发回。
                  \item \zd{复核身份审死缓，绝对不得加重刑罚}（变轻容易变重难）。救他一命，不是送他一程。
                        \begin{enumerate}
                          \item \al{直接改判、发回重审、维持原判\xhx{不再纠正} \cuo}。
                        \end{enumerate}
                \end{enumerate}
          \item 死缓限制减刑。死缓 + 限制减刑。同死缓。\zx{累犯、杀伤奸抢投毒火爆、有组织的暴力性犯罪}。
        \end{enumerate}
\end{enumerate}
